%----------------------------------------------------------------------------------------
% EVENWICHT EN GROEI. Alle opgaven zijn eigen werk.
%
% Opbouw: Probleem, Start, leerstof, en op een nieuwe pagina de Opdrachten, Koppel en Test jezelf.
% Labels van dit hoofdstuk beginnen met h2-.
%----------------------------------------------------------------------------------------

\mychapter{2}{Evenwicht en groei}

\begin{probleem}
\naast{In de cursus Matrices (hoofdstuk 4) groeide een populatie kerkuilen met ongeveer 3,2~\% per jaar. Van de jonge uilen overleefde 40~\% hun eerste jaar. Nu komt er een drukke weg door hun gebied, en sterven er meer jonge uilen. De rest van het model blijft: van klasse 2 overleeft 60~\%, en een vrouwtje krijgt gemiddeld 1,5 vrouwelijke jongen in klasse 2 en 2 in klasse 3.

Hoe groot moet de overlevingskans $s$ van de jonge uilen minstens zijn, opdat de populatie op lange termijn niet uitsterft? En hoe is een populatie die net stabiel blijft, over de klassen verdeeld?}{\begin{tikzpicture}[font=\small]
  \node[knoop] (k1) at (0,0) {1};
  \node[knoop] (k2) at (2.2,0) {2};
  \node[knoop] (k3) at (4.4,0) {3};
  \draw[boog] (k1) -- node[kans,below=1pt] {$s$} (k2);
  \draw[boog] (k2) -- node[kans,below=1pt] {0,6} (k3);
  \draw[boog,dashed] (k2) to[out=120,in=60] node[kans,above] {1,5} (k1);
  \draw[boog,dashed] (k3) to[out=110,in=80,looseness=.9] node[kans,above] {2} (k1);
\end{tikzpicture}}
\end{probleem}

\lkr{Leerplan Complementair (III-WiCo-d): LPD 30, met de wenk over matrixmodellen voor evolutie: stabilisatie is het oplossen van $AX = X$, en dat heeft niet-triviale oplossingen als 1 een eigenwaarde is. Voorkennis: hoofdstuk 1, en uit Matrices (hoofdstuk 4) overgangsmatrices, lesliematrices en het evenwicht, met dezelfde afspraak: de kolom is \textit{van}, de rij is \textit{naar}. Voor sectie 2: de limiet van $q^n$ met $\abs{q} < 1$ (Rijen). Extra's: het verband tussen het netto reproductiegetal en de eigenwaarde 1, en de dominante eigenwaarde als groeifactor op lange termijn (zonder de stelling van Perron-Frobenius). Hoge machten en nulwaarden van een derdegraadsveelterm mogen met ICT.}

In de cursus Matrices voorspelde je een evolutie met machten van een matrix: $X_n = A^n X_0$. Op lange termijn kwam het rekentoestel vaak bij een vaste toestand uit, of bij een vaste groeifactor, maar waarom, dat bleef open. In dit hoofdstuk zie je dat eigenvectoren het antwoord geven. Een evenwicht is een eigenvector bij de eigenwaarde 1, en de groeifactor op lange termijn is een eigenwaarde.

\begin{oefeningen}[Start]
\oef{h2-start-1}

\oef{h2-verken}
\end{oefeningen}

\lkr{De Verken laat zien dat het evenwicht uit Matrices een eigenvector is, en dat een begintoestand op een eigenvector alleen met een factor vermenigvuldigt. Laat leerlingen in vraag c raden wat er met een andere begintoestand gebeurt; sectie 2 geeft het antwoord.}

%========================================================================================
\section{Een evenwicht is een eigenvector}

\subsection{Stabilisatie: \texorpdfstring{$AX = X$}{AX = X}}

In de cursus Matrices was een evenwicht een toestand $X$ die niet meer verandert: $MX = X$. Lees dat met de ogen van hoofdstuk 1: $X$ is een eigenvector bij de eigenwaarde 1. Dat geeft meteen een criterium om te weten of er een evenwicht is, en een methode om het te vinden.

\begin{theorie}
\begin{stelling}{Evenwicht en eigenwaarde 1}
De vergelijking $\formule{AX = X}$ heeft een oplossing $X \neq O$ als en slechts als 1 een eigenwaarde is van $A$, dus als en slechts als

\vspace{1mm}
\qquad $\formule{\det(A - I) = 0}$

\vspace{1mm}
De oplossingen vormen dan de eigenruimte $\Eig{1}$.
\end{stelling}

\begin{bewijs}
$AX = X$ is $AX = 1 \cdot X$: een oplossing $X \neq O$ is een eigenvector bij 1. Volgens de karakteristieke vergelijking is 1 een eigenwaarde als en slechts als $\det(A - 1 \cdot I) = 0$.
\end{bewijs}
\end{theorie}

Een evenwicht is dus nooit één toestand alleen: is $X$ een evenwicht, dan ook $2X$ en $100X$. Bij een model met aandelen kies je in $\Eig{1}$ de vector waarvan de coördinaten samen 1 zijn. Bij aantallen vermenigvuldig je die met het totaal.

\voorbeeld{}{In de refter uit Matrices is $M = \km{0{,}6 & 0{,}2 \\ 0{,}4 & 0{,}8}$, met eerst warm en dan boterhammen. Bereken het evenwicht als een eigenvector.\label{vb:h2-refter}}{%
\uitleg{We gaan na dat 1 een eigenwaarde is, en lossen dan $(M - I)X = O$ op. Uit de eigenruimte kiezen we de vector met coördinaatsom 1.}
\begin{stappen}
\det(M - I) = (-0{,}4)(-0{,}2) - 0{,}2 \cdot 0{,}4 = 0 & 1 is een eigenwaarde. \\[1ex]
M - I = \km{-0{,}4 & 0{,}2 \\ 0{,}4 & -0{,}2} \iff y = 2x & Twee keer dezelfde vergelijking. \\[3ex]
\Eig{1} = \vb{(1, 2)}, \quad X = \Bigl(\tfrac13, \tfrac23\Bigr) & Deel door $1 + 2$.
\end{stappen}
Hetzelfde evenwicht als in Matrices: op lange termijn eet een derde van de leerlingen warm. Het verschil is de aanpak: nu weten we vooraf dat er een evenwicht is, en vinden we het met één stelsel.}

\begin{werkwijze}
Ga na dat $\det(A - I) = 0$. Los $(A - I)X = O$ op: dat geeft de eigenruimte $\Eig{1}$. Kies daarin de vector met coördinaatsom 1 voor de aandelen, en vermenigvuldig met het totaal voor de aantallen.
\end{werkwijze}

\nujij{Bereken het evenwicht van $\km{0{,}7 & 0{,}4 \\ 0{,}3 & 0{,}6}$.}{$M - I = \km{-0{,}3 & 0{,}4 \\ 0{,}3 & -0{,}4}$: $3x = 4y$, dus $\Eig{1} = \vb{(4, 3)}$ en $X = \bigl(\tfrac47, \tfrac37\bigr)$.}

\begin{oefeningen}
\oef{h2-even}
\end{oefeningen}

\subsection{Elke overgangsmatrix heeft een evenwicht}

In de refter was 1 een eigenwaarde. Toeval? Een overgangsmatrix heeft kolommen met som 1, omdat iedereen ergens terechtkomt. Dat is genoeg.

\begin{theorie}
\begin{stelling}{Eigenwaarde 1 van een overgangsmatrix}
Is $M$ een vierkante matrix waarvan elke kolom de som 1 heeft, dan is 1 een eigenwaarde van $M$.
\end{stelling}

\begin{bewijs}
Bekijk $M - I$. In kolom $j$ staat het element $m_{jj} - 1$ op de hoofddiagonaal en verder de elementen $m_{ij}$, dus de som van kolom $j$ is
\bewijsregel{m_{1j} + \cdots + m_{nj} - 1 = 1 - 1 = 0 .}
Tel nu alle rijen op bij de laatste rij. Dat verandert de determinant niet (Matrices, hoofdstuk 3), en de laatste rij wordt een nulrij, want op elke plaats staat een kolomsom. Ontwikkelen naar een nulrij geeft $\det(M - I) = 0$.
\end{bewijs}
\end{theorie}

Er is dus altijd een vector $X \neq O$ met $MX = X$. Dat die vector ook positieve coördinaten heeft, zodat het een echte toestand is, en dat de toestand ernaartoe gaat, bewijzen we hier niet in het algemeen. Het klopt wel als alle elementen van $M$ positief zijn: dat is een gevolg van de \textit{stelling van Perron-Frobenius}, die buiten deze cursus valt. In sectie 2 zie je waarom het gebeurt.

\voorbeeld{}{Drie streamingdiensten $P$, $Q$ en $R$ delen een markt. Elke maand blijft 80~\% van de klanten van $P$ trouw, en stapt 10~\% over naar $Q$ en 10~\% naar $R$. Het diagram toont alle overgangen. Welk deel van de markt heeft elke dienst op lange termijn?\label{vb:h2-stream}}{%
\uitleg{De overgangsmatrix lezen we kolom per kolom af. Elke kolom heeft de som 1, dus 1 is een eigenwaarde. We rijherleiden $10(M - I)$, om met gehele getallen te rekenen.}
\figuren{\begin{tikzpicture}[font=\small]
  \node[knoop] (p) at (0,0) {P};
  \node[knoop] (q) at (3,0) {Q};
  \node[knoop] (r) at (1.5,-2.5) {R};
  \draw[boog] (p) to[bend left=18] node[kans,above] {0,1} (q);
  \draw[boog] (q) to[bend left=18] node[kans,below] {0,1} (p);
  \draw[boog] (p) to[bend left=18] node[kans,pos=.45,right=1pt] {0,1} (r);
  \draw[boog] (r) to[bend left=18] node[kans,pos=.55,left=1pt] {0,2} (p);
  \draw[boog] (q) to[bend left=18] node[kans,pos=.45,right=1pt] {0,2} (r);
  \draw[boog] (r) to[bend left=18] node[kans,pos=.55,left=1pt] {0,1} (q);
  \draw[boog] (p) to[lus=180] node[kans,left] {0,8} (p);
  \draw[boog] (q) to[lus=0] node[kans,right] {0,7} (q);
  \draw[boog] (r) to[lus=270] node[kans,below] {0,7} (r);
\end{tikzpicture}}
\begin{stappenbreed}
M = \km{0{,}8 & 0{,}1 & 0{,}2 \\ 0{,}1 & 0{,}7 & 0{,}1 \\ 0{,}1 & 0{,}2 & 0{,}7} & Kolom 1: van $P$ naar $P$, $Q$ en $R$. \\[4ex]
10(M - I) = \km{-2 & 1 & 2 \\ 1 & -3 & 1 \\ 1 & 2 & -3} \to \km{1 & 0 & -\tfrac75 \\ 0 & 1 & -\tfrac45 \\ 0 & 0 & 0} & Rijherleiden; de laatste rij wordt een nulrij, zoals de stelling voorspelt. \\[4ex]
x = \tfrac75 z, \ y = \tfrac45 z\colon\ \Eig{1} = \vb{(7, 4, 5)} & Met $z = 5$. \\[1.5ex]
X = \tfrac{1}{16}(7, 4, 5) = (0{,}4375;\ 0{,}25;\ 0{,}3125) & Deel door $7 + 4 + 5 = 16$.
\end{stappenbreed}
Op lange termijn heeft $P$ ongeveer 44~\% van de markt, $Q$ een kwart en $R$ ongeveer 31~\%. Controle: $M\km{7 \\ 4 \\ 5} = \km{5{,}6 + 0{,}4 + 1 \\ 0{,}7 + 2{,}8 + 0{,}5 \\ 0{,}7 + 0{,}8 + 3{,}5} = \km{7 \\ 4 \\ 5}$.}

\begin{oefeningen}
\oef{h2-markov3}
\end{oefeningen}

\subsection{Een stabiele populatie}

Bij een lesliematrix tellen we dieren, en hebben de kolommen geen som 1: dieren sterven en worden geboren. Een evenwicht $LX = X$ is dan een \textit{stabiele populatie}: elke klasse blijft even groot. Dat kan alleen als 1 toevallig een eigenwaarde is. Wanneer gebeurt dat? Daarvoor rekenen we de karakteristieke veelterm van een lesliematrix in het algemeen uit.

\begin{theorie}
\begin{stelling}{Karakteristieke veelterm van een lesliematrix}
De lesliematrix $L = \km{v_1 & v_2 & v_3 \\ o_1 & 0 & 0 \\ 0 & o_2 & 0}$ heeft als karakteristieke veelterm

\vspace{1mm}
\qquad $\formule{p_L(\lambda) = -\lambda^3 + v_1\lambda^2 + v_2o_1\lambda + v_3o_1o_2}$
\end{stelling}

\begin{bewijs}
Met de regel van Sarrus op $L - \lambda I$: de drie producten naar rechts zijn $(v_1 - \lambda)\lambda^2$, $0$ en $v_3o_1o_2$, die naar links $0$, $0$ en $-v_2o_1\lambda$. Dus
\bewijsregel{p_L(\lambda) = (v_1 - \lambda)\lambda^2 + v_3o_1o_2 + v_2o_1\lambda .}
\end{bewijs}

\begin{stelling}{Stabiele populatie}
1 is een eigenwaarde van $L$ als en slechts als \qquad $\formule{v_1 + o_1v_2 + o_1o_2v_3 = 1}$

\vspace{1mm}
Een stabiele populatie is dan een veelvoud van \quad $\formule{X = (1,\ o_1,\ o_1o_2)}$
\end{stelling}

\begin{bewijs}
$p_L(1) = -1 + v_1 + v_2o_1 + v_3o_1o_2$ is 0 precies als de som gelijk is aan 1. Neem dan $X = (1, o_1, o_1o_2)$. De tweede en derde rij van $LX$ geven $o_1 \cdot 1 = o_1$ en $o_2 \cdot o_1$, en de eerste rij geeft $v_1 + v_2o_1 + v_3o_1o_2 = 1$. Dus $LX = X$.
\end{bewijs}
\end{theorie}

Het getal $R_0 = v_1 + o_1v_2 + o_1o_2v_3$ heeft een mooie betekenis: het is het gemiddeld aantal vrouwelijke jongen dat een pasgeboren vrouwtje in haar hele leven krijgt. Ze krijgt er $v_1$ in klasse 1; met kans $o_1$ haalt ze klasse 2 en krijgt ze er $v_2$, en met kans $o_1o_2$ haalt ze klasse 3. Is $R_0 = 1$, dan vervangt elk vrouwtje precies zichzelf. Biologen noemen $R_0$ het \textit{netto reproductiegetal}. Voor twee klassen geldt hetzelfde, met $R_0 = v_1 + o_1v_2$ en $X = (1, o_1)$.

\voorbeeld{}{Een vogelsoort leeft hoogstens twee jaar. Een vrouwtje krijgt in haar eerste jaar gemiddeld 0,5 vrouwelijke jongen, in haar tweede jaar 2. Van de jonge vogels overleeft 25~\% het eerste jaar. Is er een stabiele populatie? Zo ja, hoe ziet ze eruit?\label{vb:h2-stabiel}}{%
\uitleg{We zetten de gegevens in een lesliematrix met twee klassen, berekenen $R_0$, en lezen de stabiele verdeling af.}
\begin{stappen}
L = \km{0{,}5 & 2 \\ 0{,}25 & 0} & Vruchtbaarheid in de eerste rij, overleving eronder. \\[3ex]
R_0 = 0{,}5 + 0{,}25 \cdot 2 = 1 & Dus 1 is een eigenwaarde. \\[1ex]
X = (1;\ 0{,}25) \sim (4, 1) & Op elke vier jonge vogels één oude.
\end{stappen}
Controle: $L\kol{4 \\ 1} = \kol{2 + 2 \\ 1} = \kol{4 \\ 1}$. Een populatie van 400 jonge en 100 oude vogels blijft elk jaar precies dezelfde.}

\begin{werkwijze}
Bereken $R_0 = v_1 + o_1v_2 + o_1o_2v_3$. Is $R_0 = 1$, dan is 1 een eigenwaarde, en is elk veelvoud van $(1, o_1, o_1o_2)$ een stabiele populatie. Zoek je een onbekende parameter, los dan $R_0 = 1$ op.
\end{werkwijze}

\nujij{Bij welke overlevingskans $o_1$ is $L = \km{0{,}2 & 1{,}6 \\ o_1 & 0}$ stabiel?}{$0{,}2 + 1{,}6\,o_1 = 1$, dus $o_1 = 0{,}5$.}

\begin{oefeningen}
\oef{h2-leslie}

\oef{h2-stabiel}
\end{oefeningen}

%========================================================================================
\section{Groei op lange termijn}

\subsection{Machten op een eigenvector}

Een stabiele populatie is uitzonderlijk. Wat gebeurt er als 1 geen eigenwaarde is? Begin met een populatie die op een eigenvector ligt.

\begin{theorie}
\begin{stelling}{Machten op een eigenvector}
Is $A\vec v = \lambda\vec v$, dan is voor elk natuurlijk getal $n$

\vspace{1mm}
\qquad $\formule{A^n\vec v = \lambda^n\vec v}$
\end{stelling}

\begin{bewijs}
Voor $n = 0$ staat er $\vec v = \vec v$. Stel dat $A^n\vec v = \lambda^n\vec v$. Dan is
\bewijsregel{A^{n+1}\vec v = A\bigl(A^n\vec v\bigr) = A\bigl(\lambda^n\vec v\bigr) = \lambda^n A\vec v = \lambda^n\lambda\vec v = \lambda^{n+1}\vec v .}
Zo volgt de formule voor elke $n$ uit die voor $n - 1$: dat is een bewijs door inductie.
\end{bewijs}
\end{theorie}

Start je op een eigenvector, dan blijft de toestand dus op dezelfde rechte, en wordt hij elke stap met $\lambda$ vermenigvuldigd: een meetkundige rij van vectoren. De verdeling over de klassen verandert niet, alleen het totaal.

\voorbeeld{}{Een kever leeft twee jaar. Met $L = \km{1 & 4 \\ 0{,}5 & 0}$ start een populatie met 400 jonge en 100 oude kevers. Bereken de populatie na 10 jaar zonder $L^{10}$ uit te rekenen.\label{vb:h2-macht}}{%
\uitleg{We kijken eerst of de begintoestand een eigenvector is. Dan geeft de stelling meteen de tiende macht.}
\begin{stappen}
L\kol{400 \\ 100} = \kol{400 + 400 \\ 200} = 2\kol{400 \\ 100} & Een eigenvector bij 2. \\[3ex]
L^{10}\kol{400 \\ 100} = 2^{10}\kol{400 \\ 100} = \kol{409\,600 \\ 102\,400} & De stelling met $\lambda = 2$ en $n = 10$.
\end{stappen}
De populatie verdubbelt elk jaar, met altijd vier jonge kevers op één oude.}

\begin{oefeningen}
\oef{h2-macht}
\end{oefeningen}

\subsection{Een combinatie van eigenvectoren}

Meestal ligt de begintoestand niet op een eigenvector. Maar heeft $A$ een basis van eigenvectoren, dan is elke begintoestand een combinatie ervan, en elk stuk gedraagt zich volgens de vorige stelling.

\begin{theorie}
\begin{stelling}{Evolutie in eigenvectoren}
Zijn $\vec v_1$ en $\vec v_2$ eigenvectoren van $A$ bij $\lambda_1$ en $\lambda_2$, en is $X_0 = c_1\vec v_1 + c_2\vec v_2$, dan is

\vspace{1mm}
\qquad $\formule{X_n = A^nX_0 = c_1\lambda_1^n\vec v_1 + c_2\lambda_2^n\vec v_2}$
\end{stelling}

\begin{bewijs}
Met de rekenregels voor het matrixproduct is $A^n(c_1\vec v_1 + c_2\vec v_2) = c_1A^n\vec v_1 + c_2A^n\vec v_2$, en volgens de vorige stelling is $A^n\vec v_k = \lambda_k^n\vec v_k$.
\end{bewijs}
\end{theorie}

Met drie of meer eigenvectoren werkt het net zo, met één term per eigenvector. Het resultaat is een \textit{expliciete formule} voor $X_n$: je hoeft geen machten van matrices meer uit te rekenen.

\voorbeeld{}{Bij de vogels van voorbeeld~\ref{vb:h2-stabiel}, met $L = \km{0{,}5 & 2 \\ 0{,}25 & 0}$, start een populatie met 100 jonge en 100 oude vogels. Geef een formule voor de populatie na $n$ jaar. Wat gebeurt er op lange termijn?\label{vb:h2-combinatie}}{%
\uitleg{We hebben beide eigenwaarden en eigenvectoren nodig. Dan schrijven we $X_0$ als combinatie: dat is een stelsel voor $c_1$ en $c_2$.}
\naast{\begin{stappenvrij}{5.6cm}{3.9cm}
\lambda_1 = 1, \ \lambda_2 = -0{,}5 & Spoor $0{,}5$, determinant $-0{,}5$. \\[1ex]
\vec v_1 = (4, 1), \quad \vec v_2 = (2, -1) & $L + 0{,}5I = \km{1 & 2 \\ 0{,}25 & 0{,}5}$. \\[2ex]
4c_1 + 2c_2 = 100, \ c_1 - c_2 = 100 & $X_0 = c_1\vec v_1 + c_2\vec v_2$. \\[1ex]
c_1 = 50, \ c_2 = -50 & Tel de eerste en twee keer de tweede op. \\[1ex]
X_n = 50\kol{4 \\ 1} - 50\,(-0{,}5)^n\kol{2 \\ -1} & De stelling.
\end{stappenvrij}}{\begin{tikzpicture}[x=.42cm,y=.018cm,font=\scriptsize]
  \draw[black!14,very thin,ystep=50] (0,150) grid (8,300);
  \draw[->,grijs] (0,150) -- (8.6,150) node[right,text=black] {$n$};
  \draw[->,grijs] (0,150) -- (0,310) node[above,text=black] {totaal};
  \foreach \t in {0,2,...,8} \node[below] at (\t,150) {\t};
  \foreach \t in {150,200,250,300} \node[left] at (0,\t) {\t};
  \draw[accent!60,dashed] (0,250) -- (8,250);
  \draw[accent,thick] plot[mark=*,mark size=1.2pt] coordinates {(0,200) (1,275) (2,237.5) (3,256.25) (4,246.875) (5,251.5625) (6,249.21875) (7,250.390625) (8,249.8046875)};
\end{tikzpicture}}
Het totaal is $250 - 50(-0{,}5)^n$, want de coördinaten van $\vec v_2$ hebben som 1. Omdat $(-0{,}5)^n \to 0$, verdwijnt de tweede term: de populatie schommelt eerst, en gaat dan naar de stabiele populatie van 200 jonge en 50 oude vogels.}

\begin{werkwijze}
Bereken de eigenwaarden en eigenvectoren. Schrijf $X_0 = c_1\vec v_1 + c_2\vec v_2$ en los het stelsel op. Dan is $X_n = c_1\lambda_1^n\vec v_1 + c_2\lambda_2^n\vec v_2$. Controleer met $n = 0$ en $n = 1$.
\end{werkwijze}

\begin{oefeningen}
\oef{h2-comb}
\end{oefeningen}

\subsection{De dominante eigenwaarde}

In het voorbeeld won de eigenwaarde met de grootste absolute waarde: de term met $(-0{,}5)^n$ verdween. Dat gebeurt altijd, en het verklaart wat je in Matrices met het rekentoestel zag.

\begin{theorie}
Een eigenwaarde $\lambda_1$ van $A$ is \textbf{\textcolor{accent}{dominant}} als $\formule{\abs{\lambda_1} > \abs{\lambda}}$ voor elke andere eigenwaarde $\lambda$ van $A$ (ook de complexe).

\begin{stelling}{Gedrag op lange termijn}
Is $\lambda_1$ dominant en $X_0 = c_1\vec v_1 + c_2\vec v_2$ met $c_1 \neq 0$, dan is voor grote $n$

\vspace{1mm}
\qquad $\formule{X_n \approx c_1\lambda_1^n\vec v_1}$ \qquad want \qquad $\dfrac{X_n}{\lambda_1^n} \to c_1\vec v_1$
\end{stelling}

\begin{bewijs}
Deel de formule voor $X_n$ door $\lambda_1^n$:
\bewijsregel{\dfrac{X_n}{\lambda_1^n} = c_1\vec v_1 + c_2\Bigl(\dfrac{\lambda_2}{\lambda_1}\Bigr)^n\vec v_2 .}
Omdat $\bigl|\tfrac{\lambda_2}{\lambda_1}\bigr| < 1$, gaat $\bigl(\tfrac{\lambda_2}{\lambda_1}\bigr)^n$ naar 0 (Rijen), en dus de hele tweede term.
\end{bewijs}
\end{theorie}

Op lange termijn wordt de toestand dus elke stap met $\lambda_1$ vermenigvuldigd, en ligt hij op de rechte van $\vec v_1$. Voor een populatie betekent dat: $\lambda_1$ is de \textit{groeifactor} op lange termijn, en $\vec v_1$ geeft de \textit{stabiele leeftijdsverdeling}. Is $\lambda_1 > 1$, dan groeit de populatie; is $\lambda_1 < 1$, dan sterft ze uit. Bij een overgangsmatrix met positieve elementen is 1 de dominante eigenwaarde: alles gaat naar het evenwicht. In de refter is de andere eigenwaarde $1{,}4 - 1 = 0{,}4$ (het spoor), en $0{,}4^7 \approx 0{,}002$: daarom zag je in Matrices na een week geen verschil meer met het evenwicht. Is er geen dominante eigenwaarde, dan kan de toestand blijven schommelen: dat zie je in de opdrachten.

Voor drie eigenvectoren is het bewijs hetzelfde. Bij een lesliematrix van orde 3 is de dominante eigenwaarde een nulwaarde van een derdegraadsveelterm. Die bereken je met ICT, of door in te zoomen op de grafiek van $p_L$.

\voorbeeld{}{Bereken de groeifactor op lange termijn en de stabiele leeftijdsverdeling van de uilen uit Matrices, met $L = \km{0 & 1{,}5 & 2 \\ 0{,}4 & 0 & 0 \\ 0 & 0{,}6 & 0}$.\label{vb:h2-dominant}}{%
\uitleg{De karakteristieke veelterm volgt uit de stelling over lesliematrices. Haar nulwaarden zoeken we met ICT. Een eigenvector vinden we zoals in het bewijs van de stelling over de stabiele populatie: de overlevingskansen geven $x_2$ en $x_3$.}
\naast{\begin{stappenvrij}{6.4cm}{3.2cm}
p_L(\lambda) = -\lambda^3 + 0{,}6\lambda + 0{,}48 & $v_2o_1 = 0{,}6$ en $v_3o_1o_2 = 0{,}48$. \\[1ex]
\lambda_1 \approx 1{,}0320 & De enige reële nulwaarde. \\[1ex]
\lambda_{2,3} \approx -0{,}516 \pm 0{,}446\,i & $\abs{\lambda_{2,3}} \approx 0{,}68 < \lambda_1$. \\[1ex]
\vec v_1 = \Bigl(1,\ \dfrac{0{,}4}{\lambda_1},\ \dfrac{0{,}24}{\lambda_1^2}\Bigr) \approx (1;\ 0{,}388;\ 0{,}225) & Uit $o_1x_1 = \lambda_1x_2$ en $o_2x_2 = \lambda_1x_3$. \\[3ex]
\dfrac{\vec v_1}{1{,}613} \approx (0{,}620;\ 0{,}240;\ 0{,}140) & Deel door de som.
\end{stappenvrij}}{\begin{tikzpicture}[x=1.5cm,y=.9cm,font=\scriptsize]
  \draw[black!14,very thin,xstep=.5,ystep=.5] (-.5,-1) grid (1.5,1);
  \draw[->,grijs] (-.5,0) -- (1.6,0) node[right,text=black] {$\lambda$};
  \draw[->,grijs] (0,-1) -- (0,1.15) node[above,text=black] {$p_L$};
  \foreach \t/\l in {-.5/{-0{,}5},.5/{0{,}5},1/1} \node[below] at (\t,0) {$\l$};
  \foreach \t in {-1,1} \node[left] at (0,\t) {$\t$};
  \kromme{-.5}{1.25}{-(\x)^3+.6*\x+.48}
  \fill[fel!85!black] (1.032,0) circle (2pt);
\end{tikzpicture}}
Op lange termijn groeit de populatie dus met 3,2~\% per jaar, met 62~\% jonge uilen, 24~\% in klasse 2 en 14~\% in klasse 3. Dat zijn precies de getallen die je in Matrices met $X_{20}$ en $X_{21}$ vond, maar nu zonder hoge machten en met de reden erbij.}

\begin{werkwijze}
Bereken alle eigenwaarden, met ICT als het moet, en zoek de dominante. Dat is de groeifactor op lange termijn. Een eigenvector erbij, gedeeld door de som van zijn coördinaten, geeft de verdeling op lange termijn.
\end{werkwijze}

\begin{oefeningen}
\oef{h2-dom}
\end{oefeningen}

\subsection{Terug naar het probleem}

\voorbeeld{}{Terug naar het probleem. Hoe groot moet de overlevingskans $s$ van de jonge uilen minstens zijn opdat de populatie niet uitsterft? Hoe is een populatie verdeeld die net stabiel blijft?\label{vb:h2-terug}}{%
\uitleg{De populatie blijft op lange termijn even groot als de dominante eigenwaarde 1 is. Dat gebeurt als $R_0 = 1$. Daarna gaan we na dat 1 dan echt dominant is.}
\begin{stappenbreed}
R_0 = 0 + s \cdot 1{,}5 + s \cdot 0{,}6 \cdot 2 = 2{,}7\,s & Met $v_1 = 0$, $v_2 = 1{,}5$, $v_3 = 2$ en $o_2 = 0{,}6$. \\[1ex]
2{,}7\,s = 1 \iff s = \tfrac{10}{27} \approx 0{,}370 & De grenswaarde. \\[1.5ex]
p_L(\lambda) = -\lambda^3 + \tfrac59\lambda + \tfrac49 = -(\lambda - 1)\bigl(\lambda^2 + \lambda + \tfrac49\bigr) & Met $s = \tfrac{10}{27}$; Horner met de nulwaarde 1. \\[2ex]
\lambda_{2,3} = -\tfrac12 \pm \tfrac{\sqrt7}{6}\,i, \quad \abs{\lambda_{2,3}} = \sqrt{\tfrac14 + \tfrac{7}{36}} = \tfrac23 & 1 is dominant. \\[2ex]
X = \bigl(1,\ \tfrac{10}{27},\ \tfrac{10}{27} \cdot 0{,}6\bigr) \sim (27, 10, 6) & De stabiele populatie. \\[1ex]
\tfrac{1}{43}(27, 10, 6) \approx (0{,}628;\ 0{,}233;\ 0{,}140) & Deel door $27 + 10 + 6 = 43$.
\end{stappenbreed}
Met $s > \tfrac{10}{27}$ is $R_0 > 1$, en blijkt de dominante eigenwaarde groter dan 1 (opdracht~\ref{oef:h2-o20} bewijst dat voor twee klassen). De overlevingskans van de jonge uilen mag dus van 40~\% dalen tot ongeveer 37~\%, niet verder. Een populatie die net stabiel blijft, bestaat voor 63~\% uit jonge uilen, 23~\% in klasse 2 en 14~\% in klasse 3. Omdat de andere eigenwaarden absolute waarde $\tfrac23$ hebben, gaat elke begintoestand er snel naartoe.}

Het probleem toont de kracht van eigenwaarden: zonder één macht van $L$ uit te rekenen, vind je een grens voor een overlevingskans. Met het rekentoestel zou je moeten proberen: $s = 0{,}38$, $s = 0{,}37$, \ldots, telkens met een hoge macht.

%========================================================================================
\opdrachtenkop

\begin{oefeningen}
\oef{h2-o1}

\oef{h2-o2}

\oef{h2-o3}

\oef{h2-o4}

\oef{h2-o5}

\oef{h2-o6}

\oef{h2-o7}

\oef{h2-o8}

\oef{h2-o9}

\oef{h2-o10}

\oef{h2-o11}

\oef{h2-o12}

\oef{h2-o13}

\oef{h2-o14}

\oef{h2-o15}

\oef{h2-o16}

\oef{h2-o17}

\oef{h2-o18}

\oef{h2-o19}

\oef{h2-o20}

\oef{h2-o21}
\end{oefeningen}

%----------------------------------------------------------------------------------------
% KOPPEL: verbanden met andere leerstof
%----------------------------------------------------------------------------------------
\needspace{20\baselineskip}
\opdrachtsectie{Koppel}

\lkr{Rijen en Exponentiële en logaritmische functies zijn gezien in pakket 4. Iteraties en fractalen is een module van Complementair: geef die oefening alleen als de leerlingen ze zagen.}

\begin{oefeningen}
\oef{h2-koppel-1}

\oef{h2-koppel-2}

\oef{h2-koppel-3}
\end{oefeningen}

%----------------------------------------------------------------------------------------
\needspace{30\baselineskip} % de test en het schema Fout? Herhaal samen op een pagina
\opdrachtsectie{Test jezelf}

\begin{oefeningen}
\oef{h2-t1}

\oef{h2-t2}

\oef{h2-t3}

\oef{h2-t4}

\oef{h2-t5}

\oef{h2-t6}

\oef{h2-t7}
\end{oefeningen}

\herhaalschema{%
  \fv{oef:h2-t1}{vb:h2-refter}{oef:h2-even}%
  \fv{oef:h2-t2}{vb:h2-stream}{oef:h2-markov3}%
  \fv{oef:h2-t3}{vb:h2-stabiel}{oef:h2-stabiel}%
  \fv{oef:h2-t4}{vb:h2-macht}{oef:h2-macht}%
  \fv{oef:h2-t5}{vb:h2-combinatie}{oef:h2-comb}%
  \fv{oef:h2-t6}{vb:h2-dominant}{oef:h2-dom}%
  \fv{oef:h2-t7}{vb:h2-terug}{oef:h2-o7}}

\ifzelfstudie
\section*{Maak je eigen samenvatting}
Sluit dit boek en schrijf in je schrift: waarom een evenwicht een eigenvector bij 1 is en hoe je het berekent, waarom elke overgangsmatrix de eigenwaarde 1 heeft (met bewijs), de karakteristieke veelterm van een lesliematrix en de voorwaarde voor een stabiele populatie, wat $R_0$ betekent, de formule $A^n\vec v = \lambda^n\vec v$, hoe je $X_n$ schrijft met eigenvectoren, en wat een dominante eigenwaarde zegt over de lange termijn. Vergelijk daarna met de kaders.
\fi
